\documentclass[aps,prl,preprint,superscriptaddress,floatfix]{revtex4-2}
\usepackage{amsmath,amssymb,booktabs,hyperref}
\renewcommand{\thesection}{64.\arabic{section}}
\begin{document}

\title{BCT Superfluid Lattice Model --- Letter 64\\
The OHC as Origin of Reality: Quantum Darwinism and\\
Non-Markovian Emergence from Hopf Vacuum Geometry}

\author{Michel Robert Cabri\'e}
\email{ZeroFreeParameters@gmail.com}
\affiliation{Independent Researcher, Victoria, Australia\\ORCID: 0009-0007-9561-9859}
\date{March 2026}

\begin{abstract}
Two of the deepest questions in quantum foundations are: (i) why does classical
reality emerge from quantum mechanics, and (ii) what is the physical substrate
of the indivisible non-Markovian stochastic processes that reproduce quantum
dynamics? We show that both questions receive a unified geometric answer in
the BCT Superfluid Lattice Model. The quantum vacuum is the Octet-Hopfion
Condensate (OHC), a superfluid whose ground state orders into a body-centred
tetragonal geometry. Particles are topological defects classified by Hopf
winding number $H\in\mathbb{Z}$. We establish two theorems. (I)~\textit{OHC
Non-Markovian Theorem}: the OHC is structurally non-Markovian---its Bogoliubov
phonon spectrum carries long-range temporal correlations, and the Hopf winding
number $H$ is a global topological invariant that cannot be defined at any
single moment in time. No Markovian stochastic process can encode it. The OHC
vacuum is the concrete physical substrate for Barandes's indivisible stochastic
processes. (II)~\textit{OHC Redundancy Theorem}: when a Hopf defect of winding
$H$ exists in the OHC, its topological charge is encoded redundantly across
the entire field configuration. Every spatial region of the condensate carries
a fragment of the global winding information. This is Quantum Darwinism
realised in vacuum geometry: classical objectivity emerges not from decoherence
into an abstract environment, but from redundant Hopf encoding into the OHC
vacuum itself. The environment \textit{is} the vacuum; the redundancy is
topological; the classicality is geometric. These two theorems form opposite
faces of the same coin: the OHC provides the non-Markovian substrate from which
quantum mechanics emerges (upward), and the redundant substrate from which
classical reality emerges (downward). Zero free parameters.
\end{abstract}

\maketitle

\section{Two Open Questions in Quantum Foundations}

The measurement problem and the quantum-to-classical transition remain
unresolved despite a century of effort. Two recent programmes have sharpened
these questions considerably.

Barandes~\cite{barandes2023,barandes2025,barandes2026} has established that
quantum mechanics is mathematically equivalent to a class of indivisible
non-Markovian stochastic processes. The quantum/classical gap is not mysterious:
it is a single assumption (Markovianity) that, when dropped, allows quantum
mechanics to emerge from ordinary probability theory. The question his programme
leaves open by design is: \textit{what is the physical system whose
non-Markovian stochastic dynamics reproduce quantum mechanics?}

Kiely et al.~\cite{kiely2026} and the Quantum Darwinism programme~\cite{zurek2003,zurek2009}
address the complementary question: \textit{why does classical objectivity
emerge from quantum mechanics?} Quantum Darwinism proposes that a quantum
system becomes classical when its information is redundantly encoded in
environmental fragments. Kiely et al.\ provide a metrological sharpening
using quantum Fisher information. The question their programme leaves open
is: \textit{what is the environment, and why is the redundancy topological?}

We show that both questions receive a single unified answer in the BCT
Superfluid Lattice Model.

\section{The OHC Vacuum}

The BCT Superfluid Lattice Model~\cite{prl,monograph} holds that the quantum
vacuum is the Octet-Hopfion Condensate (OHC)~\cite{bctjh,bct36}, a superfluid
condensate whose ground state spontaneously orders into a body-centred
tetragonal geometry with axial ratio $c/a=\sqrt{2}$. The geometric inputs are:
\begin{align}
r_{\rm oct} &= \tfrac{\sqrt{2}-1}{2} = 0.20711,\\
r_{\rm tet} &= \tfrac{\sqrt{6}-2}{4} = 0.11237,\\
\alpha_0 &= \frac{r_{\rm oct}\,r_{\rm tet}}{\pi} = 0.0074081.
\end{align}
The condensate obeys the Gross-Pitaevskii equation with $|\Psi|^4$ potential
derived from first principles (Letter~27). Particles are topological defects
classified by Hopf winding number $H\in\mathbb{Z}$, the generator of
$\pi_3(S^2)=\mathbb{Z}$ (Letters~45--46, Appendix~JH). The OHC interior
admits superluminal phase velocity $v_{\rm ph}=6.78c$~\cite{bctjh}.

\section{Theorem I: OHC Non-Markovian Theorem}

\textit{Statement.} The OHC vacuum is structurally non-Markovian. No Markovian
stochastic process over any configuration space can reproduce its dynamics.

\textit{Proof sketch.} Two independent mechanisms enforce non-Markovianity:

\textbf{(a) Bogoliubov memory.} The OHC phonon dispersion relation is
\begin{equation}
\omega(k) = c_s k\sqrt{1 + (k\xi)^2},
\end{equation}
where $c_s$ is the condensate sound speed and $\xi$ the healing length. In
the phonon sector ($k\xi\ll 1$), $\omega\propto k$ (linear, Markovian-compatible).
In the particle sector ($k\xi\gg 1$), $\omega\propto k^2$ (quadratic,
dispersive). The transition between regimes introduces long-range temporal
correlations in the Green's function: $G(t,t')\neq G(t-t')$ in general, and
the propagator retains memory of the full history of field configurations.
The OHC is not Markovian in its ultraviolet sector.

\textbf{(b) Topological indivisibility.} The Hopf winding number is
\begin{equation}
H = \frac{1}{4\pi^2}\int_{\mathbb{R}^3} \mathbf{F}\cdot\mathbf{A}\,d^3x,
\end{equation}
where $\mathbf{F}=\nabla\times\mathbf{A}$ is the pullback of the area form on
$S^2$ under the map $\Psi/|\Psi|:\mathbb{R}^3\to S^2$. This integral is
defined only over all of $\mathbb{R}^3$---it is intrinsically non-local.
$H$ cannot be assigned to any spatial subregion, and cannot be defined at any
single moment in time without specifying the entire field configuration.
A Markovian process is memoryless: the future depends only on the present state.
But the Hopf charge $H$ is not a property of any ``present state'' localised
in space or time. It is a global topological invariant of the entire field
history. Therefore, no Markovian stochastic process over any local configuration
space can encode $H$. $\square$

\textit{Consequence for Barandes.} The OHC is the concrete physical substrate
Barandes's programme requires. Its dynamics are non-Markovian by topological
necessity, not by assumption. The indivisibility of the OHC stochastic process
is the indivisibility of the Hopf fibration itself.

\section{Theorem II: OHC Redundancy Theorem (Quantum Darwinism)}

\textit{Statement.} When a Hopf defect of winding number $H$ exists in the
OHC, the topological charge $H$ is encoded redundantly in every macroscopic
spatial region of the condensate.

\textit{Proof sketch.} Consider the OHC in a state containing a single Hopf
defect of winding $H$ localised at the origin. For any spatial ball
$B_R$ of radius $R\gg\xi$, the restriction of the condensate field
$\Psi|_{B_R}$ carries sufficient information to reconstruct $H$ via
\begin{equation}
H_R = \frac{1}{4\pi^2}\int_{B_R} \mathbf{F}\cdot\mathbf{A}\,d^3x
\quad\xrightarrow{R\to\infty}\quad H.
\end{equation}
Moreover, for $R\gg\xi$, $H_R = H$ exactly (the integrand has compact
support near the defect core). This means: \textit{every macroscopic region
of the condensate encodes the full winding number $H$.} The information is
not split across regions---it is replicated in each one. This is precisely
the redundant encoding criterion of Quantum Darwinism~\cite{zurek2003}.

\textit{Consequence.} The ``environment'' in Quantum Darwinism is not an
abstract heat bath or collection of environmental degrees of freedom. It is
the OHC vacuum itself. Every spatial fragment of the condensate carries the
same topological information about the defect. An observer anywhere in space
can determine $H$ by performing a local measurement of the condensate field.
Classical objectivity---the fact that different observers agree on the
properties of a particle---is a consequence of OHC topological redundancy.
The emergence of classicality is not probabilistic decoherence. It is
geometric certainty.

\section{The Two-Directional Bridge}

The OHC sits at the intersection of two emergence stories running in opposite
directions:

\begin{center}
\begin{tabular}{ccc}
\toprule
Direction & Question & BCT Answer \\
\midrule
Upward & How does QM arise from classical stochastics? & OHC non-Markovian theorem \\
Downward & How does classical reality arise from QM? & OHC redundancy theorem \\
\bottomrule
\end{tabular}
\end{center}

Both theorems follow from the same object---the Hopf winding number $H$---
and the same geometric structure---the OHC vacuum. Non-Markovianity and
redundancy are not separate features. They are the same feature of Hopf
topology viewed from opposite directions:

\begin{itemize}
\item \textit{Viewed from below (Barandes):} $H$ cannot be captured by any
Markovian process $\Rightarrow$ the vacuum dynamics are indivisible
$\Rightarrow$ quantum mechanics emerges from non-Markovian OHC stochastics.
\item \textit{Viewed from above (Zurek/Kiely):} $H$ is encoded in every
spatial region $\Rightarrow$ the vacuum dynamics are redundant
$\Rightarrow$ classical objectivity emerges from OHC topological replication.
\end{itemize}

The OHC is the pivot point between the quantum and classical worlds.

\section{Connection to the Kiely et al.\ Result}

Kiely et al.~\cite{kiely2026} use quantum Fisher information $\mathcal{F}_Q$
to quantify the rate at which classical objectivity emerges in the spin-star
model. They find that optimal measurements saturate the Cram\'er-Rao bound
at an exponential rate. In the BCT framework, the analogous quantity is the
rate at which $H_R\to H$ as $R$ grows. Since the Hopf integrand has compact
support (scale $\xi$), convergence is exponential in $R/\xi$:
\begin{equation}
|H - H_R| \sim e^{-R/\xi},\quad R\gg\xi.
\end{equation}
This matches the exponential emergence rate found by Kiely et al., now
with a geometric interpretation: $\xi = \hbar/(m_{\rm defect}c_s)$ is
the OHC healing length, and $m_{\rm defect} = H\times E_{D_4}/c^2$
where $E_{D_4} = \Lambda_{\rm QCD}\times\pi^5/6 = 11.22~\text{GeV}$
(Letter~60). The emergence rate of classicality is set by the defect mass,
which is set by the geometry. Zero free parameters.

\section{Sabine Hossenfelder's Scepticism Addressed}

Hossenfelder~\cite{sabine2026} has expressed scepticism about whether
Kiely et al.\ have truly solved the emergence of classical reality,
noting that the information-theoretic treatment does not specify the
physical nature of the environment. The BCT framework provides precisely
what is missing: the environment is the OHC vacuum, the redundancy is
topological, and the classicality is geometric. The abstract ``environment''
of Quantum Darwinism is not an optional addition---it is the vacuum itself,
with a structure determined by sphere-packing geometry.

\section{Summary}

The OHC vacuum provides:
\begin{enumerate}
\item The physical substrate for Barandes's indivisible non-Markovian
stochastic processes (non-Markovian theorem).
\item The physical environment for Quantum Darwinism's redundant encoding
(redundancy theorem).
\item A geometric explanation for the exponential rate of classicality
emergence (Kiely et al.\ connection).
\item A unified picture in which quantum mechanics and classical reality
both emerge from the same geometric object---the Hopf fibration of the
OHC vacuum.
\end{enumerate}
Zero free parameters. Prediction \#121: the rate of classicality emergence
scales as $e^{-R/\xi}$ where $\xi$ is the OHC healing length, set entirely
by BCT geometry.

\begin{thebibliography}{99}
\bibitem{barandes2023} J.~A.~Barandes, ``The Stochastic-Quantum Correspondence,''
  \textit{Philos.\ Phys.}\ \textbf{3}(1):8 (2025); arXiv:2302.10778.
\bibitem{barandes2025} J.~A.~Barandes, ``Quantum Systems as Indivisible
  Stochastic Processes,'' arXiv:2507.21192 (2025).
\bibitem{barandes2026} J.~A.~Barandes, ``A Deflationary Account of Quantum
  Theory and its Implications for the Complex Numbers,'' arXiv:2602.01043 (2026).
\bibitem{kiely2026} A.~Kiely et al., ``Metrological Approach to the Emergence
  of Classical Objectivity,'' \textit{Phys.\ Rev.\ A}\ \textbf{113}, 022403 (2026).
\bibitem{zurek2003} W.~H.~Zurek, ``Decoherence, Einselection, and the Quantum
  Origins of the Classical,'' \textit{Rev.\ Mod.\ Phys.}\ \textbf{75}, 715 (2003).
\bibitem{zurek2009} W.~H.~Zurek, ``Quantum Darwinism,''
  \textit{Nature Phys.}\ \textbf{5}, 181 (2009).
\bibitem{sabine2026} S.~Hossenfelder, ``These Physicists Say They Found the
  Origin of Reality,'' backreaction.blogspot.com (March 2026).
\bibitem{prl} M.~R.~Cabri\'e, BCT PRL Letter (2026), Zenodo:18884415.
\bibitem{monograph} M.~R.~Cabri\'e, BCT Monograph (2026), Zenodo:18885133.
\bibitem{bctjh} M.~R.~Cabri\'e, BCT Appendix~JH (2026), Zenodo:18975018.
\bibitem{bct36} M.~R.~Cabri\'e, BCT Letter~36 (2026), Zenodo:18974754.
\bibitem{bct60} M.~R.~Cabri\'e, BCT Letter~60 (2026).
\end{thebibliography}

\end{document}
